\chapter{環境構築と作業場所}

環境構築の手順は \href{../handouts/01-environment-and-workflow.md}{第1回 環境構築とワークフロー} にまとめています。
Windows・WSL・macOSの分岐，GitHub認証，VS Codeの任意設定も第1回を参照してください。

\section{教材repoと共有環境}

作業の順序は，教材repoをcloneし，そのディレクトリに移動してから，第1回で指定したPython環境を用意し，実行確認する，です。
\texttt{uv sync --locked} は教材repoの \texttt{pyproject.toml} と \texttt{uv.lock} がある場所で実行します。
clone前の空のディレクトリで実行しないでください。
依存関係は教材repoの \texttt{.venv} にそろえ，以後の回でも同じ環境を使います。

\section{演習の保存先}

第1回の動作確認用 \texttt{exercises/} は教材repoに置きます。
第2〜13回の提出物は授業で案内されたprivateな提出repoに置き，\texttt{scripts/}，\texttt{data/}，\texttt{outputs/} を使います。
提出repoのURLは授業の案内に従ってください。
具体的な実行方法と入力データのコピー先は \href{../handouts/README.md#作業場所と保存先}{handouts共通案内} を参照してください。
